\par\addvspace{12pt}\Needspace{9\baselineskip}

\section{圆锥曲线的基本性质}

% ITEM: quadratic-example-026
\solutionitem{quadratic-example-026}{例4.26}
一椭圆的中心为原点，焦点在\(x\)轴上，短轴长为\(2\)，离心率为\(\sqrt2/2\)。求椭圆方程及两条准线。若已知的是长轴、短轴的两个端点\(A(3,0)\)、\(B(0,2)\)，应怎样确定椭圆？
\solutionanswer
由\(b=1,e^2=1-b^2/a^2=1/2\)，得\(a^2=2,c=1\)。方程为\(x^2/2+y^2=1\)，准线\(x=\pm a^2/c=\pm2\)。若给定长短轴端点\((3,0),(0,2)\)，则直接读出\(a=3,b=2\)，方程为\(x^2/9+y^2/4=1\)。


% ITEM: quadratic-example-027
\solutionitem{quadratic-example-027}{例4.27}
设\(A,B\)是两定点，\(|AB|=6\)。动点\(P\)满足\(|PA|+|PB|=10\)。选择坐标系，求点\(P\)的轨迹方程及所围区域的面积，求\(|PA|\)的取值范围。
\solutionanswer
取\(AB\)中点为原点，\(A=(-3,0),B=(3,0)\)。由\(a=5,c=3\)，得\(b=4\)，所以轨迹为\(x^2/25+y^2/16=1\)，面积\(20\pi\)。焦半径在两长轴端点取到\(a-c=2,a+c=8\)，范围为\([2,8]\)。


% ITEM: quadratic-exercise-17-01
\solutionitem{quadratic-exercise-17-01}{习题4.17（1）}
分别求\(x^2/25+y^2/9=1\)与\(x^2/9+y^2/25=1\)的焦点、顶点、准线、离心率及面积。
\par
\solutionanswer
第一条椭圆中\(a=5,b=3,c=\sqrt{a^2-b^2}=4\)，焦点为\((\pm4,0)\)，顶点为\((\pm5,0),(0,\pm3)\)，离心率\(e=4/5\)，准线\(x=\pm a/e=\pm25/4\)，面积\(15\pi\)。
第二条椭圆长轴沿\(y\)轴，仍有\(a=5,b=3,c=4\)。焦点为\((0,\pm4)\)，顶点为\((0,\pm5),(\pm3,0)\)，离心率\(4/5\)，准线\(y=\pm25/4\)，面积亦为\(15\pi\)。


% ITEM: quadratic-exercise-17-02
\solutionitem{quadratic-exercise-17-02}{习题4.17（2）}
椭圆两焦点为\((\pm2,0)\)，且经过\((0,3)\)。求其方程，并求曲线上点到右焦点距离的最大值与最小值。
\par
\solutionanswer
焦点间距给出\(c=2\)。点\((0,3)\)在椭圆上，故短半轴\(b=3\)，并有\(a^2=b^2+c^2=13\)。方程为
\[\frac{x^2}{13}+\frac{y^2}{9}=1.\]
到右焦点\((2,0)\)的距离为焦半径，在右顶点\((13^{1/2},0)\)取最小值、左顶点取最大值，分别为\(a-c=\sqrt{13}-2\)与\(a+c=\sqrt{13}+2\)。


% ITEM: quadratic-exercise-17-03
\solutionitem{quadratic-exercise-17-03}{习题4.17（3）}
椭圆的一个焦点为\((0,\sqrt5)\)，一条准线为\(y=9/\sqrt5\)。已知中心为原点，求椭圆方程。
\par
\solutionanswer
长轴沿\(y\)轴，焦距参数\(c=\sqrt5\)。对应正向焦点的准线为\(y=a^2/c=9/\sqrt5\)，故\(a^2=9\)，即\(a=3\)。于是\(b^2=a^2-c^2=4\)，椭圆方程为
\[\frac{x^2}{4}+\frac{y^2}{9}=1.\]


% ITEM: quadratic-exercise-17-04
\solutionitem{quadratic-exercise-17-04}{习题4.17（4）}
设\(P\)在椭圆\(x^2/9+y^2/4=1\)上，两焦点为\(F_1,F_2\)。若\(\angle F_1PF_2=90^\circ\)，求\(\triangle PF_1F_2\)的面积。
\par
\solutionanswer
椭圆参数为\(a=3,c=\sqrt5\)，焦点间距为\(2\sqrt5\)，且对任意曲线上点\(PF_1+PF_2=2a=6\)。角\(F_1PF_2=90^\circ\)，所以以\(F_1F_2\)为斜边的直角三角形满足\(PF_1^2+PF_2^2=(2\sqrt5)^2=20\)。因此
\[PF_1\cdot PF_2=\frac{(PF_1+PF_2)^2-(PF_1^2+PF_2^2)}2=8,\]
三角形面积为\(\tfrac12PF_1PF_2=4\)。


% ITEM: quadratic-exercise-17-05
\solutionitem{quadratic-exercise-17-05}{习题4.17（5）}
（据2023年新高考I卷第5题）椭圆 \(x^2/a^2+y^2=1\)（\(a>1\)）与 \(x^2/4+y^2=1\) 的离心率分别为 \(e_1,e_2\)。若 \(e_2=\sqrt3\,e_1\)，求 \(a\)。再比较面积相同的两个椭圆：离心率较大者的长、短轴怎样变化？
\par
\solutionanswer
由 \(e_2=\sqrt3/2\) 得 \(e_1=1/2\)，所以 \(1-1/a^2=1/4\)，即 \(a=2/\sqrt3\)。一般地，椭圆的半轴比 \(b/a=\sqrt{1-e^2}\)。面积相同即 \(ab\) 相同，故离心率越大，\(b/a\) 越小，长轴越长、短轴越短。


% ITEM: quadratic-example-028
\solutionitem{quadratic-example-028}{例4.28}
双曲线的渐近线为\(y=\pm2x\)，一个焦点为\((\sqrt5,0)\)。求双曲线方程及准线方程。若只知道两条渐近线，能否唯一确定双曲线？
\solutionanswer
焦点在 \(x\) 轴上，故 \(b/a=2\)，再由 \(a^2+b^2=5\) 得 \(a=1,b=2\)。双曲线为 \(x^2-y^2/4=1\)，准线为 \(x=\pm1/\sqrt5\)。若只给两条渐近线，则整族 \(x^2-y^2/4=\lambda\)（\(\lambda\ne0\)）都符合；\(\lambda\) 的符号决定开口方向，绝对值决定尺度。


% ITEM: quadratic-example-029
\solutionitem{quadratic-example-029}{例4.29}
两定点\(F_1,F_2\)间的距离为\(10\)，动点\(P\)满足\(\bigl||PF_1|-|PF_2|\bigr|=6\)。求其轨迹方程、渐近线与离心率。若改为\(|PF_1|-|PF_2|=6\)，轨迹有何变化？
\solutionanswer
取\(F_1=(-5,0),F_2=(5,0)\)。由\(a=3,c=5\)，得\(b=4\)，轨迹为\(x^2/9-y^2/16=1\)，渐近线\(y=\pm4x/3\)，离心率\(5/3\)。若保留有符号距离差\(|PF_1|-|PF_2|=6\)，只留下右支。


% ITEM: quadratic-exercise-18-01
\solutionitem{quadratic-exercise-18-01}{习题4.18（1）}
求\(x^2/9-y^2/16=1\)的焦点、顶点、准线、离心率和渐近线，并写出其共轭双曲线方程。
\par
\solutionanswer
此双曲线有\(a=3,b=4,c=\sqrt{a^2+b^2}=5\)。焦点为\((\pm5,0)\)，顶点为\((\pm3,0)\)，离心率\(e=c/a=5/3\)，准线为\(x=\pm a/e=\pm9/5\)，渐近线为\(y=\pm(b/a)x=\pm4x/3\)。共轭双曲线方程为
\[\frac{x^2}{9}-\frac{y^2}{16}=-1,\quad\text{即}\quad\frac{y^2}{16}-\frac{x^2}{9}=1.\]


% ITEM: quadratic-exercise-18-02
\solutionitem{quadratic-exercise-18-02}{习题4.18（2）}
一双曲线的中心在原点，焦点在\(y\)轴上，离心率为\(2\)，且经过\((3,2)\)。求其方程。
\par
\solutionanswer
焦点在\(y\)轴上且中心为原点，设标准式为\(y^2/a^2-x^2/b^2=1\)，其中\(c^2=a^2+b^2\)。离心率\(c/a=2\)给出\(b^2=3a^2\)。代入经过点\((3,2)\)的条件：\(4/a^2-9/b^2=1\)，于是\(4/a^2-3/a^2=1\)，得\(a^2=1,b^2=3\)。所求方程为
\[
y^2-\frac{x^2}{3}=1.
\]
代回\((3,2)\)得\(4-3=1\)，满足条件。


% ITEM: quadratic-exercise-18-03
\solutionitem{quadratic-exercise-18-03}{习题4.18（3）}
分别讨论\(x^2-4y^2=t\)在\(t>0\)、\(t<0\)、\(t=0\)时的图形，说明这些图形与同一对直线的关系。
\par
\solutionanswer
\(t>0\)时，\(x^2-4y^2=t\)化为\(x^2/t-y^2/(t/4)=1\)，是横轴为实轴的双曲线。\(t<0\)时令\(-t>0\)，改写为\(y^2/((-t)/4)-x^2/(-t)=1\)，是纵轴为实轴的双曲线。\(t=0\)时
\[
x^2-4y^2=(x-2y)(x+2y)=0,
\]
为两条相交直线。三种情形的渐近线（退化情形本身即直线对）均为\(x=2y\)与\(x=-2y\)，也就是\( x^2-4y^2=0\)的零点直线。


% ITEM: quadratic-exercise-18-04
\solutionitem{quadratic-exercise-18-04}{习题4.18（4）}
双曲线\(x^2/4-y^2/5=1\)右支上一点到右焦点的距离为\(5\)。求该点到左焦点的距离及横坐标。
\par
\solutionanswer
右支上到左、右焦点的距离之差为 \(2a=4\)，所以到左焦点的距离是 \(9\)。两焦点为 \((\pm3,0)\)，将距离平方相减，\(12x=9^2-5^2=56\)，故 \(x=14/3\)。


% ITEM: quadratic-exercise-18-05
\solutionitem{quadratic-exercise-18-05}{习题4.18（5）}
把\(x^2/9-y^2/4=1\)化为\(uv=1\)。求\(u=v\)在原坐标中的图形，并说明能否在新坐标中直接用\(\sqrt{u^2+v^2}\)计算原点到点的距离。
\par
\solutionanswer
按题设渐近坐标\(u=x/3+y/2\)、\(v=x/3-y/2\)，原双曲线化为\(uv=1\)，逆关系为\(x=3(u+v)/2\)、\(y=u-v\)。条件\(u=v\)在原坐标中表示直线\(y=0\)；与双曲线\(uv=1\)联立还得\(u=v=\pm1\)，故交点对应原坐标\((3,0)\)、\((-3,0)\)。这些坐标轴沿渐近线方向，并非正交单位坐标。不能直接以\(\sqrt{u^2+v^2}\)计算原欧氏距离；原点到点的距离应为
\[
\sqrt{x^2+y^2}=\sqrt{\frac94(u+v)^2+(u-v)^2}.
\]


% ITEM: quadratic-exercise-18-06
\solutionitem{quadratic-exercise-18-06}{习题4.18（6）}
（2022年浙江卷第16题）过双曲线 \(x^2/a^2-y^2/b^2=1\)（\(a,b>0\)）的左焦点 \(F\) 作斜率为 \(b/(4a)\) 的直线，与双曲线交于 \(A\)，与渐近线交于 \(B\)，其中 \(x_A<0<x_B\)。若 \(|FB|=3|FA|\)，求离心率。
\par
\solutionanswer
取 \(X=x/a,\ Y=y/b\)，则双曲线为 \(X^2-Y^2=1\)，直线为 \(Y=(X+e)/4\)。因 \(X_B>0\)，在两条渐近线中只有 \(Y=X\) 给出 \(B=(e/3,e/3)\)。同一直线上的分比不变，故 \(A=F\pm(B-F)/3\)，即 \((-5e/9,e/9)\) 或 \((-13e/9,-e/9)\)。代入双曲线，分别得 \(24e^2=81\) 或 \(168e^2=81\)；后者与 \(e>1\) 矛盾，所以 \(e=3\sqrt6/4\)。


% ITEM: quadratic-example-030
\solutionitem{quadratic-example-030}{例4.30}
（2025年全国II卷第6题）设抛物线\(C:y^2=2px\)（\(p>0\)）的焦点为\(F\)，点\(A\)在\(C\)上，过\(A\)作准线的垂线，垂足为\(B\)。若直线\(BF\)的方程为\(y=-2x+2\)，求\(|AF|\)。
\solutionanswer
焦点\((p/2,0)\)在\(BF\)上，故\(0=-p+2\)，得\(p=2\)。准线为\(x=-1\)，与\(BF\)交于\(B=(-1,4)\)。\(AB\)垂直准线，故\(A=(4,4)\)。由焦准线定义，\(|AF|=|AB|=5\)。


% ITEM: quadratic-exercise-19-01
\solutionitem{quadratic-exercise-19-01}{习题4.19（1）}
分别写出\(y^2=8x\)、\(x^2=-4y\)的顶点、焦点、准线与对称轴。
\par
\solutionanswer
\(y^2=8x\)对应\(2p=8\)，故顶点为\((0,0)\)、焦点为\((2,0)\)、准线为\(x=-2\)，对称轴为\(x\)轴。\(x^2=-4y\)对应\(x^2=-2py\)且\(p=2\)，故顶点为\((0,0)\)、焦点为\((0,-1)\)、准线为\(y=1\)，对称轴为\(y\)轴。


% ITEM: quadratic-exercise-19-02
\solutionitem{quadratic-exercise-19-02}{习题4.19（2）}
抛物线顶点为\((1,-2)\)，焦点为\((3,-2)\)。求其方程和准线。
\par
\solutionanswer
顶点到焦点的距离为\(2\)。本节标准式\((y-k)^2=2p(x-h)\)中焦点为\((h+p/2,k)\)，故\(p/2=2\)、\(p=4\)，方程为
\[
(y+2)^2=8(x-1).
\]
准线为\(x=h-p/2=-1\)。焦点到准线距离为\(4\)，与标准参数\(p=4\)一致。


% ITEM: quadratic-exercise-19-03
\solutionitem{quadratic-exercise-19-03}{习题4.19（3）}
（2024年上海卷第7题）抛物线\(y^2=4x\)上一点到准线的距离为\(9\)，求该点到\(x\)轴的距离。
\par
\solutionanswer
准线为\(x=-1\)，故\(x+1=9\)，\(x=8\)。代入得\(y^2=32\)，到\(x\)轴距离为\(4\sqrt2\)。


% ITEM: quadratic-exercise-19-04
\solutionitem{quadratic-exercise-19-04}{习题4.19（4）}
一拱形截面为抛物线，拱顶高出水面\(2\)米，水面处宽\(4\)米。若水位下降\(1\)米，求新的水面宽度，并说明所选坐标系的作用。
\par
\solutionanswer
取水面为\(x\)轴、拱顶正下方为原点，抛物线为\(y=2-x^2/2\)。新水面\(y=-1\)与之相交于\(x=\pm\sqrt6\)，宽为\(2\sqrt6\)米。利用对称轴建系，使一次项消失，高度也直接由纵坐标表示。


% ITEM: quadratic-example-031
\solutionitem{quadratic-example-031}{例4.31}
（1）（2012年安徽卷文科第14题）过抛物线\(y^2=4x\)焦点\(F\)的直线交抛物线于\(A,B\)。若\(|AF|=3\)，求\(|BF|\)。
\par



（2）过椭圆\(x^2/9+y^2/4=1\)右焦点的弦\(AB\)满足\(|AF|=2|BF|\)。求两段焦半径及弦长，并判断这样的弦是否存在。
\par
\solutionanswer
（1）抛物线参数为\(p=2\)，焦弦满足\(1/|AF|+1/|BF|=2/p=1\)，所以\(|BF|=3/2\)。

（2）椭圆焦弦满足\(1/|AF|+1/|BF|=2a/b^2=3/2\)。令两距离为\(2s,s\)，得\(s=1\)，两焦半径为\(2,1\)，弦长为\(3\)。这样的弦确实存在：右焦半径取遍 \([3-\sqrt5,3+\sqrt5]\)，其中含 \(2\)。取焦半径为 \(2\) 的端点作焦弦，另一段由倒数和必为 \(1\)。


% ITEM: quadratic-example-032
\solutionitem{quadratic-example-032}{例4.32}
设\(F_1,F_2\)是椭圆\(x^2/a^2+y^2/b^2=1\)的两焦点，\(P\)是椭圆上一点。求\(|PF_1|\,|PF_2|\)的取值范围；若\(\angle F_1PF_2=90^\circ\)，求三角形面积，并讨论这种位置的点何时存在。
\solutionanswer
取两焦点为\((\pm c,0)\)。由焦半径\(a\pm ex\)，乘积为\(a^2-e^2x^2\)，故范围为\([b^2,a^2]\)。直角时两焦半径平方和为\(4c^2\)，而和为\(2a\)，所以乘积为\(2b^2\)，面积为\(b^2\)。存在这样的点当且仅当\(2b^2\le a^2\)，即\(e\ge1/\sqrt2\)；等号时点在短轴端点，仍为非退化直角三角形。


% ITEM: quadratic-exercise-20-01
\solutionitem{quadratic-exercise-20-01}{习题4.20（1）}
求\(x^2/16+y^2/7=1\)、\(x^2/9-y^2/16=1\)、\(y^2=6x\)的通径端点与通径长。
\par
\solutionanswer
椭圆通径端点为\((\pm3,\pm7/4)\)，两焦点各对应一条竖直通径，长度均\(7/2\)。双曲线通径端点为\((\pm5,\pm16/3)\)，长度均\(32/3\)。抛物线通径端点为\((3/2,\pm3)\)，长度为\(6\)。


% ITEM: quadratic-exercise-20-02
\solutionitem{quadratic-exercise-20-02}{习题4.20（2）}
过\(y^2=4x\)焦点的弦长为\(8\)。求两段焦半径的长度及弦所在直线的方程。
\par
\solutionanswer
设两焦半径为 \(r,s\)。由 \(r+s=8,\ 1/r+1/s=1\)，得 \(rs=8,\ |r-s|=4\sqrt2\)，两段长为 \(4\pm2\sqrt2\)。准线为 \(x=-1\)，故每个端点的焦半径就是其横坐标加 \(1\)。弦在 \(x\) 轴上的投影长因此为 \(|r-s|=4\sqrt2\)，恰是弦长的 \(1/\sqrt2\)。所求两条焦弦直线为 \(y=\pm(x-1)\)。


% ITEM: quadratic-exercise-20-03
\solutionitem{quadratic-exercise-20-03}{习题4.20（3）}
设\(P\)是椭圆\(x^2/25+y^2/9=1\)上一点。若\(|PF_1|=2|PF_2|\)，求点\(P\)的坐标和\(\triangle PF_1F_2\)的面积。
\par
\solutionanswer
取\(F_1=(-4,0),F_2=(4,0)\)。距离和为\(10\)，所以两焦半径为\(20/3,10/3\)。由\(|PF_1|=5+4x/5\)，得\(x=25/12\)，再代入得\(y=\pm\sqrt{119}/4\)。面积为\(\frac12\cdot8\cdot|y|=\sqrt{119}\)。交换焦点编号时横坐标相应取反。


% ITEM: quadratic-exercise-20-04
\solutionitem{quadratic-exercise-20-04}{习题4.20（4）}
对抛物线的一条焦弦\(AB\)，以\(AB\)为直径作圆。证明此圆与准线相切，并说明圆心到准线的距离与半径的关系。
\par
\solutionanswer
设圆心为弦中点\(M\)。焦点在抛物线的两交点之间，故半径\(R=(|AF|+|BF|)/2\)。两端点到准线的距离分别就是焦半径，又位于准线同侧，所以\(M\)到准线的距离为两距离的平均值，也等于\(R\)。因此圆与准线相切。


% ITEM: quadratic-exercise-20-05
\solutionitem{quadratic-exercise-20-05}{习题4.20（5）}
对双曲线\(x^2/9-y^2/16=1\)，分别考察过右焦点的直线\(y=0\)与\(x=5\)截出的弦，判断焦点与弦的位置关系，说明上述倒数和公式能否直接用于双曲线。
\par
\solutionanswer
\(y=0\)时端点为\((-3,0),(3,0)\)，右焦点\((5,0)\)在弦外，两距离为\(8,2\)。\(x=5\)时端点为\((5,\pm16/3)\)，焦点在弦中点，两距离均为\(16/3\)。前一情形倒数和为\(5/8\)，后一情形为\(3/8\)，可见椭圆式的普通距离倒数和公式不能直接无条件用于双曲线，必须注意焦点与弦的位置和有向关系。


% ITEM: quadratic-exercise-20-06
\solutionitem{quadratic-exercise-20-06}{习题4.20（6）}
设\(A(x_1,y_1),B(x_2,y_2)\)是抛物线\(y^2=2px\)（\(p>0\)）的一条焦弦的端点。证明\(y_1y_2=-p^2\)，并求\(x_1x_2\)。
\par
\solutionanswer
焦弦可写为\(x=p/2+ky\)，代入得\(y^2-2pky-p^2=0\)，故\(y_1y_2=-p^2\)。由\(x_i=y_i^2/(2p)\)，进一步得到\(x_1x_2=p^2/4\)。


% ITEM: quadratic-exercise-20-07
\solutionitem{quadratic-exercise-20-07}{习题4.20（7）}
过椭圆 \(x^2/4+y^2=1\) 的左焦点 \(F_1\) 作弦 \(AB\)，右焦点为 \(F_2\)。求 \(\triangle AF_2B\) 的最大面积及取等时弦的方向。
\par
\solutionanswer
\(F_1\) 在弦上，以 \(F_1F_2\) 为公共底边，两小三角形的面积之和是 \(\sqrt3\,|y_A-y_B|\)。既然只需纵坐标之差，就把非水平弦写成 \(x=my-\sqrt3\)。代入椭圆得 \((m^2+4)y^2-2\sqrt3\,my-1=0\)。不必分别求根，直接用根差得到
\[
 S=\frac{4\sqrt3\sqrt{1+m^2}}{m^2+4}
   =\frac{4\sqrt3\,w}{w^2+3}\le2,\qquad w=\sqrt{1+m^2}.
\]
当且仅当 \(w=\sqrt3\)，即 \(m=\pm\sqrt2\) 时取等；此时弦的斜率为 \(\pm1/\sqrt2\)。水平弦的面积为零。


% ITEM: quadratic-example-033
\solutionitem{quadratic-example-033}{例4.33}
定长为\(L\)的线段\(AB\)两端分别沿互相垂直的两条直线滑动，点\(P\)在线段上，且\(AP=tL\)，其中\(0<t<1\)固定。求点\(P\)的轨迹，并说明\(t=1/2\)时有什么特点。
\solutionanswer
取滑动直线为两坐标轴，\(A=(u,0),B=(0,v)\)，则\(u^2+v^2=L^2\)。\(P=(1-t)A+tB\)，故
\[
\frac{x^2}{(1-t)^2L^2}+\frac{y^2}{t^2L^2}=1.
\]
反向令\(u=x/(1-t),v=y/t\)即可恢复线段，故这就是完整轨迹。\(t=1/2\)时为半径\(L/2\)的圆。


% ITEM: quadratic-exercise-21-01
\solutionitem{quadratic-exercise-21-01}{习题4.21（1）}
判断\(x^2/9+y^2/4=1\)与\(x^2/18+y^2/8=1\)是否相似，写出一个将前者变为后者的变换。再与\(x^2/9+y^2=1\)比较。
\par
\solutionanswer
前两椭圆相似，作\((x,y)\mapsto(\sqrt2x,\sqrt2y)\)即可。第三条的两半轴比为\(3\)，与前两条的\(3/2\)不同，所以不相似。


% ITEM: quadratic-exercise-21-02
\solutionitem{quadratic-exercise-21-02}{习题4.21（2）}
写出把\(y^2=4x\)变为\(y^2=12x\)的一次相似变换。若要求变换后的方程为\(y^2=12x+8\)，应再作怎样的平移？
\par
\solutionanswer
先作相似变换\((X,Y)=(3x,3y)\)，得\(Y^2=12X\)。再沿\((-2/3,0)\)平移，即得\(Y^2=12X+8\)。


% ITEM: quadratic-exercise-21-03
\solutionitem{quadratic-exercise-21-03}{习题4.21（3）}
半径为\(4\)的圆经上述正投影后得到离心率为\(\sqrt3/2\)的椭圆。求倾角、椭圆面积，并解释面积倍率。
\par
\solutionanswer
正投影椭圆的半轴为\(4,4\cos\theta\)，故\(e=\sin\theta\)。题设给出\(\theta=60^\circ\)，半轴为\(4,2\)，面积\(8\pi\)。投影沿一方向不变，另一方向缩为\(\cos\theta=1/2\)，面积倍率也为\(1/2\)。


% ITEM: quadratic-exercise-21-04
\solutionitem{quadratic-exercise-21-04}{习题4.21（4）}
证明一条双曲线与其共轭双曲线相似，当且仅当它是等轴双曲线。
\par
\solutionanswer
原双曲线实半轴、共轭半轴为\(a,b\)，共轭曲线的相应长度为\(b,a\)。若相似比为\(k\)，必须\(ka=b,kb=a\)，从而\(a=b\)。反之，\(a=b\)时旋转\(90^\circ\)就将两曲线互换，故条件充分必要。


% ITEM: quadratic-exercise-21-05
\solutionitem{quadratic-exercise-21-05}{习题4.21（5）}
长为\(5\)的线段\(AB\)两端分别沿\(x\)轴、\(y\)轴滑动，点\(P\)在线段上且\(AP=2\)。求\(P\)的轨迹；若改为\(AP=3\)，所得两条轨迹是否相似？写出一个将它们对应起来的变换。
\par
\solutionanswer
写\(A=(u,0),B=(0,v)\)，则\(u^2+v^2=25\)。\(AP=2\)时\(P=(3u/5,2v/5)\)，所以轨迹为\(x^2/9+y^2/4=1\)。\(AP=3\)时成为\(x^2/4+y^2/9=1\)。旋转\(90^\circ\)或关于\(y=x\)镜像，均可使两轨迹对应，因此相似。


% ITEM: quadratic-example-034
\solutionitem{quadratic-example-034}{例4.34 *}
将椭圆\(x^2/9+y^2/4=1\)化为单位圆。分别求两条直线\(y=x\)、\(y=-x\)的像，判断变换前后是否垂直；再对\(y=x\)、\(y=-4x/9\)作同样研究，并联系共轭直径解释结果。
\solutionanswer
取\(u=x/3,v=y/2\)。第一对直线成为\(v=\pm3u/2\)，不再垂直；第二对成为\(v=3u/2,v=-2u/3\)，反而垂直。第二对原直径满足斜率乘积\(-4/9\)，是椭圆的共轭直径，映到圆后成为垂直直径。这体现了共轭关系的保持与普通垂直关系的不保持。


% ITEM: quadratic-exercise-22-01
\solutionitem{quadratic-exercise-22-01}{习题4.22（1） *}
写出将\(x^2/9+y^2/4=1\)变为单位圆的矩阵，求两焦点的像，并说明这两个点在像圆中的位置。
\par
\solutionanswer
取\(U=\operatorname{diag}(1/3,1/2)X\)。焦点\((\pm\sqrt5,0)\)映为\((\pm\sqrt5/3,0)\)，都是单位圆内部的点，既不是圆心，也不是像圆的焦点。仿射变换不保持焦点这种度量性质。


% ITEM: quadratic-exercise-22-02
\solutionitem{quadratic-exercise-22-02}{习题4.22（2） *}
设\(q(X)=X^THX\)正定，\(H=S^{-T}S^{-1}\)。证明变换\(X=SU\)将\(q(X)=1\)及其内部变为单位圆及其内部，并将\(H\)的双线性正交关系变为普通内积的正交关系。
\par
\solutionanswer
由\(S^THS=I\)，有\(q(SU)=U^TU\)。因此边界\(q=1\)与单位圆对应，内部\(q<1\)与单位圆盘对应。又\(b(SU,SV)=U^TV\)，所以由\(H\)决定的正交关系正好变为普通坐标内积的正交关系。


% ITEM: quadratic-exercise-22-03
\solutionitem{quadratic-exercise-22-03}{习题4.22（3） *}
证明可逆线性变换\(X=SU\)保持所有向量的长度，当且仅当\(S^TS=I\)；若所有长度都乘以固定正数\(t\)，相应条件是什么？
\par
\solutionanswer
保长等价于\(U^T(S^TS-I)U=0\)对所有\(U\)成立。括号内为对称矩阵，取标准基向量及它们的和，便知各元素均为零，故\(S^TS=I\)。反向直接代入即可。若长度统一乘\(t>0\)，同理条件为\(S^TS=t^2I\)。
\par
